% from ega4-1-fr.tex
\oldpage[IV-1]{222}
\thispagestyle{plain}

\begin{center}
{\Large CHAPITRE IV}

\vspace{1.5cm}
{\LARGE\bfseries ÉTUDE LOCALE DES SCHÉMAS\\[0.3em]
ET DES MORPHISMES DE SCHÉMAS}

\vspace{1.5cm}
{\large\bfseries Sommaire}
\end{center}

\medskip
\begin{tabular}{@{}r p{0.84\textwidth}}
§ 1. & Conditions de finitude relatives. Ensembles constructibles dans les préschémas.\\
§ 2. & Changement de base et platitude.\\
§ 3. & Cycles premiers associés et décompositions primaires.\\
§ 4. & Changement du corps de base dans les préschémas.\\
§ 5. & Dimension et profondeur dans les préschémas.\\
§ 6. & Morphismes plats de préschémas localement noethériens.\\
§ 7. & Application aux relations entre un anneau local noethérien et son complété.
Anneaux excellents.\\
§ 8. & Limites projectives de préschémas.\\
§ 9. & Propriétés constructibles.\\
§ 10. & Préschémas de Jacobson.\\
§ 11.\footnotemark[1] & Propriétés topologiques des morphismes plats de présentation finie.
Critères locaux de platitude.\\
§ 12. & Étude des fibres des morphismes plats de présentation finie.\\
§ 13. & Morphismes équidimensionnels.\\
§ 14. & Morphismes universellement ouverts.\\
§ 15. & Étude des fibres d'un morphisme universellement ouvert.\\
§ 16. & Invariants différentiels. Morphismes différentiellement lisses.\\
§ 17. & Morphismes lisses, non ramifiés, étales.\\
§ 18. & Compléments sur les morphismes étales. Anneaux locaux henséliens.\\
§ 19. & Immersions régulières et transversalement régulières.\\
§ 20. & Sections hyperplanes~; projections génériques.\\
§ 21. & Prolongements infinitésimaux.\\
\end{tabular}
\footnotetext[1]{L'ordre et le contenu des §§ 11 à 21 ne sont donnés qu'à titre
indicatif et pourront être modifiés avant leur publication.}

\clearpage
